\fichatitulo{53 · SÍNTESE DE MÉTRICAS}{arXiv · 2024}{Avaliação de RAG: survey}
{\small Yu et al.. \textit{Evaluation of Retrieval-Augmented Generation: A Survey}. arXiv · 2024. \href{https://arxiv.org/abs/2405.07437}{Fonte consultada}.\par}
\campo{Problema e objetivo}
Como organizar métricas, conjuntos de dados e limitações da avaliação de sistemas RAG?
\campo{Principal contribuição · síntese interpretativa}
Oferece um mapa unificado das dimensões de recuperação e geração, incluindo relevância, correção e fidelidade.
\campo{Método / natureza da evidência}
Survey comparativo de benchmarks e métricas, com proposta de um processo unificado de avaliação.
\campo{Resultado ou argumento principal}
A estrutura híbrida do RAG exige avaliar recuperação e geração de forma articulada; benchmarks existentes cobrem dimensões e tarefas de modo desigual.
\campo{Excertos-chave · seleção literal}
\begin{itemize}
\excerto{1}{we conduct A Unified Evaluation Process of RAG}{Resumo.}
\end{itemize}
\campo{Limitações e análise crítica}
Como survey, depende da qualidade e do recorte dos estudos sintetizados; deve orientar a taxonomia, sem substituir benchmarks primários.
\campo{Aproveitamento na dissertação · proposta}
Organizar o quadro comparativo da revisão e explicitar a separação entre recuperação, geração, fidelidade e avaliação humana.
\registro{Fonte consultada nesta preparação: preprint; consulta focal do resumo e da taxonomia. Chave: \texttt{yuEtAl2024ragEvaluationSurvey}.}
\clearpage
